\cotefile{67}
\scanpage{86}
% archive sheet number: 88

qui opère sur $\mathrm{Oc}(V)$, opère sur $I(V)$, et on trouve

\[
\underset{\substack{\shortparallel\\\deleted{\unsure{G|V}}\operatorname{Aut}^{+}(\mathrm{Oc}(V))}}{\Gamma(V)}
\xrightarrow{\sim}\unsure{\mathfrak{S}|\mathscr{S}}_{I(V)}\,,
\]

et cet iso. \deleted{\illegible{1}} définit une structure de $V$-torseur sur $I(V)$. (Ceci résulte de la théorie combinatoire du $\unsure{N|V}$ octaèdre, ou du cube c'est kif kif.)

\uline{Proposition} \unsure{Soit|S\textsuperscript{t}} $X$ droite projective \unsure{rélles|réelle} sur un corps alg clos $k$ (il suffit qu'il \deleted{\unsure{contienne}} soit stable par extraction des racines carrées) de car. $\ne 2$, et $G=\operatorname{Aut}_{k}X$. Alors il y a corr 1-1 entre

a) les sous-groupes $V$ de $G$ isom : $\mathbb{D}_2$

b) \rule{5em}{0.4pt} $\Gamma$ de $G$ isom : $\mathfrak{S}_4$ \deleted{\unsure{non|un}}

c) les parties $S\in\mathcal{P}_6(\deleted{\unsure{S}}X(k))$ qui sont ``octaèdrales'', i.e. \deleted{telles que :} \inserted{isom :} $(\mathbb P^1_k,\{0,\infty,\pm1,\pm i\})$

par
\[
\begin{aligned}
\Gamma&=\operatorname{Nor}_{G}(V)\\
V&=\inserted{\text{unique }}\text{sous-groupe }\inserted{\text{inv. de }\Gamma,}\ \deleted{\unsure{\text{inv. de }G|\text{inv. de }\Gamma}}\ \text{isom : }\mathbb D_2\ \deleted{\unsure{\text{à }\mathbb B}}\\
S&=\text{ens de ramification de l'action de }V\text{ sur }X\\
&\hspace{2em}\deleted{\unsure{\text{ens de l'action}}},
\end{aligned}
\]

% slanted
\margin{\uline{Car} la catégorie \unsure{(au lieu précisé)|(au sens précisé)} des couples $(X,V)$ (ou $(X,\Gamma)$, $(X,S)$) $\to$ \uline{équivalente} : celle des octaèdres \uline{orientés}, une fois choisie une racine carrée de $-1$ dans $k$. Mieux : la cat des couples $(X,V)$ \unsure{est $\to$|éq $\to$} celle des couples $(\mathrm{Oc},\varphi)$ où $\mathrm{Oc}$ est un \unsure{polyèdre|octaèdre} : $\varphi:\{\pm i\}$ $\underset{=}{\sim}\,\unsure{\mathrm{or}|\omega}(\mathrm{Oc})$}

(\uline{NB} $\Gamma$ a un ens de ramification plus grand qui comprend, en plus de $S$, les ``milieux d'arêtes'' et ``centres des faces'' de l'octaèdre,

Dans une telle partie $S$, on a un antipodisme $\underline{a}_S=\underline{a}_V$ défini par la condition \deleted{\unsure{que}} pour $\sigma\in V^*$, $\underline{a}_S$ \deleted{soit} \inserted{induise} la transposition de $X^\sigma$ (de cardinal 2). On définit sur $S$ une structure octaèdrale, admettant $\Gamma$ comme groupe des autom orientés, et $\underline a$ comme antipodisme, via l'antipodisme. (les faces sont les parties $S'\in\mathcal{P}_3(S)$ telles que $S'\cup\underline a S'=S'$).
